\section{Relative lifting with additive rank}
\label{lifting:section}

We prove Proposition~\ref{counting:relative-target}, with returned rank
at most $d_0+d$.  The old patch is removed one weight layer at a time.
After a recursive call, we discard the old test and transfer only the
returned patch.  Removing $D$ old slots requires at most $D$ new lift
slots, even though the preparation may introduce more ambient coordinates.

Two statements supply the return from sampled paths to the parent.
The weighted major decomposition recovers an ambient slow-plus-rational
decomposition from decompositions on many paths.  Globalization then
replaces the local polynomial orbits by one ambient orbit in the same
group, preserving a specified projection.  In the patch application,
this replacement adds no slots.  Positive bump selection and
Property~\ref{sampling:property-transfer} of
Proposition~\ref{sampling:constrained} perform the remaining transfer.
Theorem~\ref{lifting:relative-theorem} assembles these steps and also
proves the impossibility alternative.  Throughout, we retain the
sampler's order of parameters and its uniform complexity bounds.

\subsection{The weighted major decomposition}

We first recover an ambient decomposition from decompositions on sampled
paths.  The conclusion separates a polynomial with bounded real
coefficients from one with controlled rational denominator, using
different variables for the two estimates.  The following parameter
convention makes this recovery uniform in the later sampling choices.

In the remainder, an \emph{early budget} $p$ bounds dimensions,
logarithms of inverse chart widths, logarithmic accuracies, log inverse
productive probability, and log inverse score.  A bound is \emph{early} if it is at most
$\mathcal P(p)$ in logarithmic complexity, or $\mathcal B(p)$ in
size, and is independent of the later sampling parameters and the
heights of the spaces $W_h$.  These later parameters must be chosen
only after all early input and output budgets have been fixed.
The sampler dimension and its structural block counts are fixed from
the input dimensions before this numerical budget, as in
Proposition~\ref{sampling:constrained}; they do not depend on chart
widths or requested accuracies.  Degree bounds depend only on the
fixed degrees.  The early numerical bounds may depend on the chart
widths, but the rank accounting below uses the actual numbers of
patch slots and the inequality $\sum_h\dim W_h\le D$.

\paragraph{Length thresholds for local symbols.}
In each local symbol argument, fix a threshold $J_{\rm loc}$ whose
logarithm has an early bound.  An axis of the current reindexed slice
is short if its side length is less than $J_{\rm loc}$; every other
axis is long.  Fixing a short axis costs at most $\log J_{\rm loc}$,
so fixing all short axes has early cost.  Homogeneous parts in long
variables always use all axes retained after these specializations.
In particular, an inactive axis with side at least $J_{\rm loc}$
remains a variable of the symbol.  These long variables need not
coincide with the parallel variables $t_\parallel$ in the forecasting
law of Section~\ref{cube:transfer-section}.
The bounded number of threshold increases used below can be
forecast from early data before the later sampling choices.

\begin{proposition}[Weighted major decomposition]
\label{lifting:candidate-major}
Use the data $U,C_h,W_h,c_h,y_h,\beta_h$ of
Proposition~\ref{sampling:constrained}, with rank at least $R$ in
Definition~\ref{sampling:rank}, and assume
Properties~\ref{sampling:property-geometry} and~\ref{sampling:property-detection} of that
proposition on the sampled boxes and the ambient domain $U$.
Assume also the uniform symbol step-drop statement of
Lemma~\ref{prelim:step-drop}.
For every fixed degree bound $s$ there is an exponent $C_s$
with the following uniform property.  Prescribe an early budget $p$,
sampling dimensions, chart widths, a fraction at least $e^{-p}$ of
productive paths, and slow and denominator bounds at most $e^p$.
The output bounds are at most
$\exp((2+p)^{C_s})$ before choosing the perturbation size, scale
range, rank threshold, or parent side lengths.  Choose those later
parameters sufficiently in that order, as in
Proposition~\ref{sampling:constrained}.

Write $A_h=(C_h)_{[h]}$ for the ordinary degree-$h$ homogeneous part
of $C_h$, and put $A=(A_h)_h$.  Give $u$ weight $1$ and $b_h$ weight
$h$.  Let $D(u,b)$ be any fixed scalar weighted homogeneous polynomial
of degree $1\le d\le s$, with no bound on its coefficients.
If, after fixing short axes on each productive path, the degree-$d$
homogeneous part of $D(\psi(t),\beta(\psi(t)))$ in all the retained
parameter variables is slow plus rational of
the prescribed bounds on at least the prescribed fraction of paths,
then there exist rational graded spaces $K_h\supset W_h$ of
early height and polynomials $S,R$ such that, for $b\in\bigoplus_hK_h$,
\begin{equation}
 D(u,b)=S(u/H,b-A(u))+R(u,b).
 \label{lifting:major-form}
\end{equation}
The coefficients of $S$ and a common denominator of $R$ have early
bounds; the numerators of $R$ are unrestricted.  The assertion is simultaneous
for polynomially many scalar polynomials, by intersecting their
spaces.  Long-variable thresholds and slice costs also have
early bounds.

\end{proposition}

To prove Proposition~\ref{lifting:candidate-major}, we will realize the
phase $e(D(u,\beta(u)))$ as an ordinary filtered niltest and apply
Lemma~\ref{prelim:step-drop}.  We first construct the group in which
this realization takes values.

\subsubsection{The translation group and its potential}

Fix an integer $d\geq1$ and finite-dimensional real vector spaces $E_h$
equipped with rational bases, $1\leq h\leq d$.
Write $E=\bigoplus_{h=1}^dE_h$, assigning weight $h$ to a coordinate of
$E_h$.  Let $\mathscr Q_{<d}(E)$ be the vector space of scalar polynomials
in $b\in E$ of weighted degree less than $d$.  A coefficient of a
monomial of weight $k$ is assigned group degree $d-k$.

The multiplication below is suggested by polynomial differences.
Let $V:E\to\R$ be weighted homogeneous of degree $d$, and for fixed
$x\in E$ put $F_x(b)=V(b)-V(b-x)$.  Its terms of weight $d$ in $b$
cancel, so $F_x\in\mathscr Q_{<d}(E)$, and for $x,z\in E$ we have
\[
 F_x(b-z)+F_z(b)=V(b)-V(b-x-z)=F_{x+z}(b).
\]
Thus translating the polynomial coordinate as in the next lemma makes
these differences compose with their base translations.  Under its
top-frequency annihilation hypothesis, Lemma~\ref{lifting:potential}
will show that the relevant subgroup elements have this difference
form on their projected space.

\begin{lemma}[Polynomial translation group]
\label{lifting:translation-group}
On $\mathscr G_d=E\times\mathscr Q_{<d}(E)$, the operation
\begin{equation}
 (x,F)(z,Q)=(x+z,F(b-z)+Q(b))
 \label{lifting:group-law}
\end{equation}
defines a connected simply connected graded nilpotent Lie group of
degree at most $d$.  Its Lie bracket is
\begin{equation}
 [(X,P),(Z,Q)]=(0,\partial_XQ-\partial_ZP).
 \label{lifting:bracket}
\end{equation}
In rational coordinate bases, integer $x$ and integer polynomial
coefficients form a lattice.  For fixed $d$, all group formulas have
polynomially many coordinates and bounded degree in the dimension of
$E$.
\end{lemma}

\begin{proof}
Both associations of a product of three elements have polynomial
coordinate
$F(b-z-w)+Q(b-w)+R(b)$.
The identity is $(0,0)$ and the inverse is
$(-x,-F(b+x))$.  Thus the operation is a group operation on a real
vector space, with polynomial multiplication and inversion.
Expanding a commutator to second order gives
\eqref{lifting:bracket}.  Differentiation in $E_h$ lowers polynomial
weight by $h$ and hence raises group degree by $h$.
Every bracket of total degree greater than $d$ vanishes.  The maps
which multiply $x_h$ by $t^h$ and the coefficient of $b^\alpha$ by
$t^{d-\mathrm{wt}(\alpha)}$ are group dilations, proving the grading.
The underlying manifold is a vector space, so it is connected and
simply connected.

Integer translations preserve integer polynomial coefficients.  This
proves closure and discreteness of the stated subgroup.  Reduce the
$x$ coordinates modulo integers and then the coefficients of $F$
modulo integers to obtain a bounded fundamental region.  Consequently
the subgroup is cocompact.  The number of monomials of ordinary degree
at most $d$ in $\dim E$ variables is at most $(1+\dim E)^d$ up to a
constant depending on $d$; this also bounds the number and degrees of
the formulas just used.
\end{proof}

The next Lemma turns top-frequency annihilation into a polynomial
potential.  A rational height bound controls both numerators and
denominators; a denominator bound alone does not control numerators.

\begin{lemma}[Potential associated with a top-annihilated algebra]
\label{lifting:potential}
Let $\mathfrak v$ be a graded rational Lie subalgebra of the Lie algebra
of $\mathscr G_d$, and let $K_h\subset E_h$ be its projection in degree
$h$.  Put $K=\bigoplus_hK_h$.  Suppose a linear functional on the top
layer vanishes on $\mathfrak v_d$ and takes the value $1$ on the
constant-polynomial element $(0,1)$.  Then there is a rational weighted
homogeneous polynomial $V:K\to\R$ of weight $d$ such that
\begin{equation}
 F(b)=V(b)-V(b-x)\qquad
 ((x,F)\in\exp\mathfrak v,\ b\in K).
 \label{lifting:potential-identity}
\end{equation}
For fixed $d$, a polynomial bound on dimensions and on logarithmic
rational heights of $\mathfrak v$ gives such a bound for $V$ and for a
rational polynomial extension of $V$ to $E$.
\end{lemma}

\begin{proof}
First let $(0,Q)\in\mathfrak v_j$ be homogeneous.  We claim
$Q|_K=0$.  Otherwise a nonzero coefficient of $Q|_K$ can be extracted
by successively differentiating in homogeneous coordinate directions
of $K$, with total weight $d-j$.  Lift each such direction to an
element of the corresponding layer of $\mathfrak v$.
By \eqref{lifting:bracket}, the iterated bracket with $(0,Q)$ is the
same iterated directional derivative, viewed as a pure-polynomial
element.  It would be a nonzero constant in $\mathfrak v_d$,
contradicting the hypothesis.  This proves the claim.

Choose rational graded linear lifts $(X,P_X)\in\mathfrak v$ of
$X\in K$.  The claim shows that $P_X|_K$ is independent of the lift.
Define a polynomial one-form by
$\omega_b(X)=P_X(b)$ for $b,X\in K$.
The bracket of two chosen lifts has zero base component, so the claim
and \eqref{lifting:bracket} give
$\partial_XP_Z=\partial_ZP_X$ on $K$.  Thus $\omega$ is closed.
Its weighted degree, including the weight of the differential, is
$d$.  If $\mathcal E_b=\sum_h h b_h$ is the weighted Euler vector
field, set
\[
                  V(b)=d^{-1}\omega_b(\mathcal E_b).
\]
The closedness identity and Euler's identity imply
$\partial_X V(b)=\omega_b(X)$.
This can also be checked on each coordinate monomial, so it requires
no choice of a primitive or integration constant.

The set of pairs with $x\in K$ satisfying
\eqref{lifting:potential-identity} is a closed subgroup: substitution
in \eqref{lifting:group-law} makes the
two differences telescope.  Its Lie algebra consists of pairs
$(X,P)$ with $X\in K$ and $P|_K=\partial_XV$.  We have proved that this Lie algebra
contains $\mathfrak v$.  Since $\exp\mathfrak v$ is connected, the
subgroup identity follows.

Rational sections, restrictions, and extensions can be found by
solving linear systems of polynomial size.  The formula for $V$ adds
only multiplication by a coordinate and division by the fixed integer
$d$.  Clearing denominators and bounding minors therefore gives the
asserted rational-height bounds.
\end{proof}

\begin{lemma}[The algebraic major decomposition]
\label{lifting:algebraic-major}
Let $D(u,b)$ be weighted homogeneous of degree $d$, and let $D_0(b)$
be its part independent of $u$.  Let $A(u)$ be a weighted homogeneous
map into $K$.  Suppose the symbol
\[
 \bigl(A(u),D(u,b)-D_0(b-A(u))\bigr)
       =(e,E)(x,F)(a,Q)
\]
has middle factor in the subgroup in
Lemma~\ref{lifting:potential}, whose defining algebra has bounded
rational height, and $e,x,a$ take values in $K$.
Suppose the left factor has coefficients bounded in normalized
variables $u_i/H_i$, and the right factor has rational coefficients
of bounded common denominator.  If $D_0$ has bounded coefficients,
then, on $b\in K$,
\[
 D(u,b)=S(u/H,b-A(u))+R(u,b),
\]
where $S$ has bounded coefficients and $R$ has bounded common
denominator.  Bounds have polynomial logarithmic dependence on those
in the hypotheses, for fixed degree.
\end{lemma}

\begin{proof}
The base identity is $A=e+x+a$.  Use
Lemma~\ref{lifting:potential} and the group law twice to obtain
\begin{align*}
 D(u,b)={}&D_0(b-A)+E(u,b-A+e)-V(b-A+e)\\
           &\quad+Q(u,b)+V(b-a).
\end{align*}
The first three terms are polynomials with bounded coefficients in
$u/H$ and $b-A$.  The last two have rational coefficients of bounded
common denominator.  This proves \eqref{lifting:major-form}.
Products and substitutions involve a bounded degree and polynomially
many coefficients, giving the stated bounds.
\end{proof}

\subsubsection{Proof of the weighted major decomposition}

\begin{proof}[Proof of Proposition~\ref{lifting:candidate-major}]
We apply Lemma~\ref{prelim:step-drop} to ordinary polynomial maps
on integer boxes.  The polynomial translation group converts the
weighted polynomial in the conclusion into this setting.

First subtract the integer parts of all coefficients of $D$.
This leaves its exponential unchanged at integer $u,b$ and changes
each leading path evaluation by an integer-coefficient polynomial.
It therefore preserves the local major hypothesis.  The number of
coefficients is polynomial in the dimensions, since $d$ is fixed.
We can hence assume they all lie in $[0,1)$, and restore the removed
integer polynomial to $R$ at the end.

Use the buffered site input
$\chi(u)e(D(u,\beta(u)))$.  On a productive path its top part can be
made nearly constant on controlled subboxes and residue classes.
Indeed a common residue modulus clears the rational coefficients of
the top part, and subdivision in normalized coordinates makes the
slow part oscillate by at most $1/100$.  Fix only the axes shorter
than the early threshold $J_{\rm loc}$.  The other inactive axes
remain in the symbol.  Since the full substituted polynomial has
ordinary degree at most $d$, the pure degree-$d$ part in the
retained variables is independent of the values fixed on short
axes.  These specializations have early slice cost.  Cancel the
remaining lower-degree terms by an ordinary polynomial phase of
degree at most $d-1$.  Choose $\chi$ to equal one on all path supports and
to vanish outside the buffered charts.  The local correlation is
then at least $9/10$.  Property~\ref{sampling:property-detection} of
the sampler, applied with degree
$d-1$, gives an early positive ambient correlation with a site twist
times a degree-$(d-1)$ niltest.

Partition normalized spatial coordinates and their bounded-modulus
residues so that the smooth spatial factors can be frozen to the
accuracy required by this correlation.  Choosing one part preserves
a correlation of early size.  At every tag $h>d$, Fourier expand the
smooth lift dependence on the ambient torus cover of the bounded
modulus used by the twist.  This is the ambient coordinate torus,
not a torus basis of $W_h$; no height of $W_h$ enters its Fourier
cutoff.  Convolution with a trigonometric approximate identity
approximates a Lipschitz function uniformly.  The coefficients remain
Lipschitz functions of the unexpanded coordinates, with the same
bound up to the fixed truncation cost.  Thus the number of terms,
coefficient mass, and remaining Lipschitz bounds are early bounded.
A residue character of $\beta_h$ is handled on the same cover by
writing $\beta_h=C_h-c_h-y_h$.  This gives a bounded rational
frequency on $C_h-c_h$ and a smooth factor of $y_h$.

Set $E_h=\R^{d_h}$ for $1\le h\le s$ and write $D_0(b)=D(0,b)$.
Since $D$ uses only tags at most $d$, we represent these factors on
$\mathscr G_d$ formed from $\bigoplus_{h\le d}E_h$.  In this
realization all base vectors are restricted to those tags.  Set
\[
 x_h=C_h(u)-c_h,\qquad
 F(u,b)=D(u,b)-D_0(b-x(u)).
\]
Its integral-part evaluation is
$e(F(u,\beta)+D_0(\beta-x))=e(D(u,\beta))$.
Right multiplication by an integral lattice element changes the
phase argument by an integer: if $(z,Q)$ is integral, then the new
integral part is $\beta+z$ and the argument increases by
$Q(\beta+z)\in\Z$.  To include residues of $\beta$ modulo $M$,
use the subgroup lattice with $z\in M\Z^{\dim E}$ and integral
$Q$; its index and grid bounds are early.  On this cover every remaining low-tag
twist is a Lipschitz function of $y=x-\beta$ and $\beta\bmod M$.
Buffered cutoffs therefore give a well-defined nilmanifold function,
including the twist, and remove the jumps where the integral part
changes.
On bounded group representatives the possible integral parts are
bounded.  The coefficient reduction therefore makes this function
Lipschitz with early bounds.  The coefficient of a dummy monomial of
weight $k<d$ in $F$ has ordinary $u$-degree at most $d-k$, while the
base $x_h$ has degree at most $h$.  The group law and the logarithmic
coordinate change respect these assigned degrees, so the same degree
bounds hold in logarithmic coordinates.  Hence this is an ordinary
filtered polynomial map of degree at most $d$.  The dummy terms of
weight $d$ cancel by the definition of $D_0$.  Taking the homogeneous
part in each group grade replaces $x_h$ by $A_h$ and gives exactly
the symbol
\[
             \bigl(A(u),D(u,b)-D_0(b-A(u))\bigr).
\]
Thus no lower-degree coefficient of $C_h$ enters this symbol.

Consider a resulting term with a highest tag $h>d$ whose frequency
does not annihilate $W_h$.  All factors at higher tags are constant
on $C_l-c_l\in W_l$.  The degree-$h$ factor is a single torus phase;
every other factor is an ordinary niltest of degree at most $h-1$.
For the factors of tags at most $d$, use the nilmanifold
realization just constructed.  A correlation larger than
an arbitrarily prescribed early error would, by
Lemma~\ref{prelim:step-drop} applied to the product group, force the
top of that torus phase to be slow plus bounded-denominator rational:
the partner has no top coordinate and the fast algebra is annihilated
by the torus frequency.  Clearing its early bounded denominator gives
an integer row forbidden by high rank.  Choose the rank threshold
large enough for all these rows and the required error.  The total
contribution of these terms is negligible, because their total
coefficient mass is early bounded.  At least one term involving only
tags at most $d$ consequently retains early correlation.  All errors
and cutoffs used in this paragraph were determined before that rank
threshold was chosen.

A top-frequency decomposition retains frequency $1$ on the constant
polynomial direction: every Fourier term of an already equivariant
function retains that equivariance.  Select a correlated term and
apply Lemma~\ref{prelim:step-drop} jointly with the degree-$(d-1)$
partner.  The partner has zero grade $d$, including when $d=1$ and
it is constant.  Consequently the selected top functional on the
product descends to a functional on the top of $\mathscr G_d$,
takes value $1$ on $(0,1)$, and annihilates the projection of the
fast algebra.  Rational projection therefore supplies precisely the
algebra required by Lemma~\ref{lifting:potential}.  Its base
projection contains $W_h$: otherwise a bounded-height rational row
annihilating that projection, but not $W_h$, would exhibit a forbidden
slow-plus-rational top component of $C_h$.  Once this inclusion is
known, write the base splitting as $A=e+x+a$ with $x\in K$.
Since $A\in W\subseteq K$, each coefficient of $e+a$ lies in $K$.
All these symbol coefficients have positive degree, so
Lemma~\ref{prelim:separation} puts the slow and rational base factors
individually in $K$.
Lemma~\ref{lifting:algebraic-major} now applies.  Undoing the spatial
residue and subbox parametrization multiplies homogeneous coefficient
bounds and denominators by early factors only.  More explicitly,
on $u=u_0+qv$ the degree-$h$ symbol is $A_h(qv)$;
if the new sides are $T_i\asymp\delta_iH_i/q$, undoing this
substitution multiplies slow bounds by at most
$\prod_i\delta_i^{-\alpha_i}$ and denominators by $q^h$.
Here $q$ and $\delta_i^{-1}$ are early bounded.  This also justifies
each use of the original rank hypothesis after restricting to a
spatial residue subbox.  For $h>d$ take
$K_h=E_h$.  For polynomially many simultaneous scalar components,
apply the argument separately with a common early budget, intersect
the resulting spaces, and restrict their identities.  A joint
stacked-row computation bounds the rational height of this
intersection by an early quantity.
\end{proof}

\begin{remark}[Uniformity of the bounds]
\label{lifting:major-gap}
The preceding proof uses Property~\ref{sampling:property-detection}
of Proposition~\ref{sampling:constrained} and
Lemma~\ref{prelim:step-drop} with their stated early bounds.
Its Fourier cutoffs use ambient lattice covers of
bounded modulus, so this step introduces no additional dependence on
the height of $W_h$.  Choosing the local model only after allowing
its complexity to depend on that height or the final long scale
would, however, invalidate this deduction.
\end{remark}

\subsection{Globalization in a fixed group}

\begin{proposition}[Globalization in a fixed group]
\label{lifting:candidate-globalization}
Assume Properties~\ref{sampling:property-geometry} and
\ref{sampling:property-detection} of Proposition~\ref{sampling:constrained},
including its rank hypothesis and order of later parameter choices,
the uniform symbol step-drop and separation statements, and their
stated polynomial logarithmic bounds.  Fix an early budget $p$ and
a strong rational projection $\pi:G^0\to Y$ of degree $r\le s$.
Include the dimensions, rational structure constants, filtration
bases, lattice grids, and matrix of $\pi$ in the early complexity
budget.  Let $P_Y(u,b)$ be a weighted filtered polynomial, with no
coefficient bound.  Let $\Phi_u:G^0/\Gamma^0\to[0,1]$ have an
early uniform Lipschitz bound.  Assume that, at each requested early
accuracy $\eta$, the family $\{\Phi_u\}$ has a supremum-norm
$\eta$-net of at most $\exp(\mathcal P(p+\log(1/\eta)))$ functions.
Its bins may be arbitrary site sets.  Let $a(u)$ be a fixed real
input of early bounded supremum.

There are early output bounds, independent of every later sampling
parameter, with the following consequence.  Suppose a fraction at
least $e^{-p}$ of paths have a permitted slice and a filtered
polynomial $g_\psi$ satisfying
\[
 \pi(g_\psi(t))=P_Y(\psi(t),\beta(\psi(t))),\qquad
 \E_t a(\psi(t))\Phi_{\psi(t)}(g_\psi(t)\Gamma^0)\ge e^{-p}.
\]
Then there is one weighted filtered polynomial $g^U(u,b)$ in the
same group, satisfying the formal identity $\pi g^U=P_Y$, such that
an early fraction of the original productive paths have permitted
subslices on which replacement by $g^U$ gives early positive score.
No coefficient bound on $g_\psi$ or $g^U$ is required.
\end{proposition}

\begin{proof}
We induct on $r$.  For $r=0$ the group is trivial and its sole map
is already fixed.  All early budgets below are computed before the
later sampling choices.  The number of successive polynomial
increases of those budgets will be bounded in terms of $s$.

\emph{Top modes and fixed ambient partners.}
Write $\delta\ge e^{-p}$ for the input score and
$B\ge\max(1,\|a\|_\infty)$ for its early bound.  Truncate the Fourier
expansion in the top torus $G^0_r$ to uniform error at most
$\delta/(100B)$, using Lemma~\ref{prelim:vertical}.  Call a mode
significant on a productive path if its absolute tested average is
at least the score divided by one hundred times the number of
truncated modes.  The sum of nonsignificant terms is then negligible.
Let $L\subset G^0_r\cap\ker\pi$ be the joint connected annihilator
of the restrictions of significant modes.  Averaging the actual
positive test over $L$ retains those modes, preserves the range
$[0,1]$, and retains early positive score.  Choose a basis among
these significant restricted frequencies.  There are at most
$\dim(G^0_r\cap\ker\pi)$ pivots.

Choose these bases in a fixed order.  If there are $M$ truncated
frequencies and the kernel top layer has dimension $D$, the number
of possible bases, and hence of possible $L$, is at most $(1+M)^D$.
Pigeonhole $L$ and the pivot basis first.  This loses only an early
fraction, rather than counting all subsets of the $M$ modes.
If there are no pivots, the partner-selection steps below are empty;
the product has no partner factors and its fast algebra is the full
graded algebra.

For each pivot approximate the external function family in supremum
norm and partition sites by a nearest bin.  One bin detects the
pivot with early correlation.  For every possible bin apply the
sampler's linear-combination approximation to $a$ times the bin
indicator.  Such indicators are allowed because its input is an
arbitrary bounded function.  Use its testing seminorm over all
permitted slices and all current local degree-$r$ niltests.
The test at this point is only the vertical component of a fixed
bin function evaluated at $g_\psi$; its ambient partner has not yet
been selected.  Its degree is $r$, and its group and function
complexity are early bounded.  Polynomial coefficients of the orbit
do not enter niltest complexity.  Conjugation and products, represented
on product groups, preserve the degree bound.
Choose the error smaller than the pivot threshold times the
productive fraction divided by the total early number of bins and
modes.  Markov's inequality and a union bound make all these
approximations valid on all but a negligible fraction of paths,
uniformly over their separately chosen slices and tests.
Explicitly, when there are pivots, let $J\ge1$ be the number of masked inputs, let $\kappa$
be the retained productive probability, and let $\tau$ be smaller
than half the least correlation threshold after selecting a bin.
Normalize each masked input by $B$ before applying
Property~\ref{sampling:property-detection}, and
choose its approximation so the unnormalized error has $X$-norm
at most $\varepsilon=\kappa\tau/(4J)$.  If $e_j(\psi)$ denotes the
pathwise supremum error for mask $j$, then
\[
 \Pr_\psi\{\max_j e_j(\psi)>\tau\}
 \le \sum_j\frac{\E_\psi e_j(\psi)}{\tau}
 \le\kappa/4.
\]
The supremum in Property~\ref{sampling:property-detection} covers all pivot frequencies as well as
the separately chosen slices.  Its input test complexity is fixed
before any model partner is selected.  Net centers can be chosen
from the original family, at twice the requested accuracy, so they
retain the common Lipschitz and range bounds.

For a retained path and pivot, some model term gives biased pairing
of the local vertical function with a site twist $J$ times a fixed
ambient partner map $Q_\ell(u)$.  Fourier truncate the top of that
partner and retain one mode.  By subbox and residue subdivision,
freeze $J$ to error smaller than the pairing; one cell retains biased
pairing without it.  Freeze axes shorter than the early threshold
required by Lemma~\ref{prelim:step-drop}.  Applying that Lemma to the
ordinary local pair gives a splitting of its pure-long symbol into
a rational algebra killed by the joint top frequency.

Refer each homogeneous identity back to the original long variables.
Translation does not alter it and residue dilation costs only early
factors.  Thus pieces used for distinct pivots need not intersect:
their symbol splittings can be joined by separation.  All long sides
exceed the fixed degree and interpolation cutoffs.  Pigeonhole the
partner choices, frequencies and bounded-height rational algebras,
which have an early number of possibilities jointly for the
polynomially many pivots.  An early fraction of productive paths now
has fixed choices.  Set
\[
 G=(G^0/L)\times\prod_\ell G_{Q_\ell},\qquad
 F=Y\times\prod_\ell G_{Q_\ell},\qquad
 P_F=(P_Y,(Q_\ell)_\ell).
\]
The product has only polynomially many factors, irrespective of the
length of the model lists.  First project each pair splitting through
$G^0\to G^0/L$; its top constraint descends because the pivot
frequency annihilates $L$.  Pull each resulting fast algebra back
to this product, leaving the unused partner coordinates free.
Each such inverse-image algebra admits a splitting of the full
local symbol: extend its outer factors by identities in those unused
coordinates.  Lemma~\ref{prelim:symbol-splitting} gives a splitting
in their intersection; this is the joined fast algebra $\mathfrak v$.
A top vector with zero projection to $F$ has zero partner
coordinates.  Its pivot constraints force the original coordinate
into $\log L$, so it is zero after quotienting.  Consequently
$\mathfrak v_r\to(\pi\mathfrak v)_r$ is injective.

\emph{Matching the fixed marked symbol.}
Write $Z$ for the weighted symbol of $P_F$ and
$\mathfrak v_F=\pi\mathfrak v$.  Initially take $K_h=E_h$ and
$\mathcal E=\mathcal R=1$.  Inductively suppose through grade $d-1$
that the logarithm of $\mathcal E^{-1}Z\mathcal R^{-1}$ belongs to
$\mathfrak v_F$ on $K=\bigoplus_hK_h$.  The left factor has early
bounded polynomial coefficients in $u/H,b-A(u)$; the right factor
has rational polynomial coefficients of early common denominator.

Raise the long threshold so that the leading evaluations of
$\beta_h$ lie in $K_h$.  Indeed their projections off $K_h$ have
rational coefficients of early denominator, since $\beta_h$ has
integer coefficients and $K_h$ has early rational height.  They are
also negatives of coefficients of the fully slow $y_h$, since
$C_h-c_h\in W_h\subset K_h$.  Separation makes each pure-long
coefficient zero.  The evaluations of the two fixed outer factors
are therefore slow and rational ordinary symbols.  Here
$\beta_h^{\rm top}-A_h(\psi^{\rm lin})=-y_h^{\rm top}$ is slow,
whereas $\psi^{\rm lin}$ and $\beta_h^{\rm top}$ are integral.

Compare the projected local splitting with these fixed factors.
Through grade $d-1$, separation puts the ratios of their corresponding
outer factors in $\mathfrak v_F$.  Lift the ratios within that
algebra and absorb them into the local fast factor.  Thus the
projected local outer factors agree exactly, in those grades, with
the evaluations of the fixed ones.  The grade-$d$ logarithmic
coordinate of the fixed residual, modulo $(\mathfrak v_F)_d$,
then evaluates to a sum of slow and rational scalar polynomials:
all lower discrepancies are zero and the fast coordinate vanishes
in the quotient.

Apply Proposition~\ref{lifting:candidate-major} simultaneously to
those scalar coordinates and intersect its spaces with $K_h$.
Rational linear sections lift the two quotient terms to $s_d,r_d$.
Update
\[
 \mathcal E\longleftarrow\mathcal E\exp(s_d),\qquad
 \mathcal R\longleftarrow\exp(r_d)\mathcal R.
\]
At grade $d$ the residual loses exactly $s_d+r_d$ and lower grades
are unchanged.  This proves the invariant in the next grade.
BCH preserves the two allowed types.  There are at most $s$ such
threshold increases, and restriction to fewer long axes preserves
all earlier homogeneous identities.

At the last threshold freeze nonlong axes on favorable fibers of
productive slices.  The whole projection of $\beta_h$ off $K_h$
is now constant: each remaining nonconstant coefficient is both
rational of early denominator and fully slow.  The constant has
early bounded norm because it is the projection of $-y_h$, and early
denominator because $\beta_h$ is integral.  Pigeonhole these finitely
many constants on an early fraction of paths.  Let $p_h$ be the
rational affine projection onto the resulting fixed translate of
$K_h$ and temporarily replace $P_F$ by
$P_F^*(u,b)=P_F(u,p(b))$.  These maps agree on the retained slices.

\emph{Full factors with exact projection.}
Refilter by Lemma~\ref{lifting:refiltration}.  Lift
$\mathcal E,\mathcal R$ to full polynomial factors $E_F,R_F$ of
$P_F^*$.  Lift their logarithmic layer formulas, replacing $b-A(u)$
on the left by $\operatorname{proj}_K b-(C-c)(u)$, and using the
rational formulas in $u,\operatorname{proj}_K b$ on the right.
Thus $E_F$ is fully slow and bounded on the sampling charts and
$R_F$ has early common denominator.  Symbol matching puts every
positive-degree coefficient of
$P_H=E_F^{-1}P_F^*R_F^{-1}$ in the appropriate $H_{F,i}$: parts in
deeper original layers already belong there because
$F_{i+1}\subset H_{F,i}$.  In the vector quotient $F/H_{F,1}$ only
a constant remains.  Split it in rational coordinates into a bounded
fractional part and an integral part, and insert lifts on the inside
of $E_F$ and $R_F$.  The middle map now takes values in $H_F$.

For each local map construct
\[
 g=EhR,\qquad \pi(E)=E_F,\quad\pi(R)=R_F,\quad\pi(h)=P_H.
\]
Here is the exact kernel correction.  Write the local symbol as
$epr$, with $p\in\exp\mathfrak v$, and write $e_F,r_F$ for the
symbols of the prescribed projected outer factors.  By
Lemma~\ref{prelim:symbol-splitting},
\[
 A=(\pi e)^{-1}e_F\in\exp\mathfrak v_F,
 \qquad D=r_F(\pi r)^{-1}\in\exp\mathfrak v_F.
\]
The first is slow and the second rational.  Lift their logarithms
through a controlled graded section to $\widehat A,\widehat D$ in
$\exp\mathfrak v$ and replace the splitting by
\[
 e'=e\widehat A,\qquad
 p'=\widehat A^{-1}p\widehat D^{-1},\qquad
 r'=\widehat D r.
\]
It now has $\pi e'=e_F$ and $\pi r'=r_F$.  Choose a filtered
rational logarithmic section $\rho$ of $\pi$ and put
$E_0=\rho(E_F)$, $R_0=\rho(R_F)$.  The symbols
$e'(\operatorname{Sym}E_0)^{-1}$ and
$(\operatorname{Sym}R_0)^{-1}r'$ lie in
$\ker(\operatorname{gr}\pi)$.  Strongness supplies controlled
actual-kernel lifts of each graded coefficient: if
$X\in\log G_i$ represents one and $\pi X\in\log F_{i+1}$,
subtract a section lift of $\pi X$ in $\log G_{i+1}$.
Lift the two symbol discrepancies coefficientwise in this way to
kernel polynomial maps $K_E,K_R$, with no other logarithmic terms.
Then $E=K_EE_0$ and $R=R_0K_R$ have the prescribed projections
exactly, retain the slow and rational bounds, and give a middle map
whose positive-degree coefficients are adapted to $H_i$.

For its constant term use the vector quotient $q:G\to G/H_1$.
Since $\pi(H)=H_F$, the condition $\pi(h(0))\in H_F$ implies
$q(h(0))\in q(\ker\pi)$.  Choose a controlled rational linear
section of $q$ on this rational subspace, split $q(h(0))$ in a
rational basis into a bounded fractional vector and an integral
vector, and lift both into $\ker\pi$.  Absorb their exponentials
on the inside of $E$ and $R$, respectively.  These constant
adjustments preserve the symbol and the exact projections; the
middle constant now lies in $H_1$.  Thus $h\in H$ and all three
projection identities hold in the groups.

Fix the logarithmic section $\rho$.  The kernel map
$E\rho(E_F)^{-1}$ is fully slow and bounded.  Early subdivision
makes it close to a fixed bounded element $k_E$ from an early net.
The rational kernel map $\rho(R_F)^{-1}R$ has early common
denominator; residue restriction makes its right lattice coset
constant, represented by a fixed bounded rational $k_R\in\ker\pi$.
If these common denominators vary with the path, first pigeonhole
their values, of which there are an early number.  A cover clears
only the finitely many denominators actually selected, rather than
every integer below an early bound.
Pigeonholing retains early score on an early fraction of paths.  Put
\[
              E^U=k_E\rho(E_F),\qquad R^U=\rho(R_F)k_R.
\]
Right lattice changes preserve cosets.  Left freezing changes the
test by at most its Lipschitz constant times the freezing error,
uniformly in the unbounded middle map, by right invariance of the
upstairs metric.

\emph{Lattice reconstruction and external function nets.}
By Lemma~\ref{lifting:refiltration}, the map
$i:H\to H/H_r\times H_F$ is injective onto the group fiber product.
We explain why finer lattices make the same assertion usable for
nilmanifold functions.  Choose a source lattice $\Gamma_H$ permitting
the right translates $R^U$.  This uses only their common denominator:
the adjoint formulas for rational elements have bounded denominators
independently of the sizes of their numerators, so a common deep
integral grid lies in all required conjugate lattices.

In rational logarithmic coordinates let $T$ be a rational left
inverse of $i_*$.  If $l\Z^{\dim H}\subset\log\Gamma_H$, choose
a sufficiently divisible early integer $D$ so that
$T(D\Z^{\dim(H/H_r)+\dim H_F})\subset l\Z^{\dim H}$.
Take deep target lattices $\Delta'',\Delta_F$ whose logarithms lie
in that target grid and set
$\Delta_H=i^{-1}(\Delta''\times\Delta_F)$.
Rationality implies that $\Delta_H$ is a lattice, and the grid
choice gives $\Delta_H\subset\Gamma_H$.  The original test pulls
back to $H/\Delta_H$, which embeds in the product of the two covered
nilmanifolds.  Rational subgroup reconstruction gives an early
Lipschitz bound on the joint image.  The quantitative reason is that
on bounded representatives any target lattice element outside the
rational image has distance at least an early inverse size from it;
nearby pairs can therefore be compared using a lattice element in
the image and the bounded rational inverse $T$.

For clarity, fix these lattices for the whole family and write $J$
for the resulting joint image.  If $\phi$ is an original function,
$e$ a bounded value of $E^U$, and $r$ a bounded representative of a
right-lattice residue of $R^U$, the function on $J$ is
\[
 f_{\phi,e,r}(i(h))=\phi(ehr\Gamma).
\]
All these functions have one early Lipschitz bound $L_*$, by
Lemma~\ref{prelim:reconstruction}; enlarge $L_*$ to be at least one
and to dominate the Lipschitz bounds of the original functions
$\phi$ as well.  For a fixed $r$, right invariance
also gives
\[
 \|f_{\phi,e,r}-f_{\phi',e',r}\|_\infty
 \le \|\phi-\phi'\|_\infty+
       \operatorname{Lip}(\phi')d_G(e,e').
\]
Use the clipped McShane extension from the same $J$ with the same
constant $L_*$ for every function.  The infimum formulas show that
this operation is nonexpansive in supremum norm.  Specialize the
extension at the second coordinate $y=P_H(u,\beta(u))\Delta_F$;
changing $y$ changes the resulting function on $(H/H_r)/\Delta''$
by at most $L_*d(y,y')$ in supremum norm.
Thus, writing $\Psi_{\phi,e,r,y}$ for this specialized function,
\[
 \|\Psi_{\phi,e,r,y}-\Psi_{\phi',e',r,y'}\|_\infty
 \le \|\phi-\phi'\|_\infty
       +L_*d_G(e,e')+L_*d(y,y').
\]

There are an early number of possible right-lattice residues $r$:
Mal'cev reduction of rational elements of one fixed common
denominator leaves bounded representatives of controlled
denominator.  This count is independent of all numerators.  The
bounded $e$-region and the compact $y$-nilmanifold have polynomial
dimension and early coordinate bounds.  Choose an original function
net at accuracy $\eta/6$ and recenter each nonempty bin at a member
of the original family, giving accuracy $\eta/3$ with its original
range and Lipschitz bounds.  Together with grids of mesh
$\eta/(3L_*)$ in $e$ and $y$ and the finite choices of $r$, this
gives a net with
logarithmic size polynomial in $p+\log(1/\eta)$.  The dependence of
$e,r,y$ on the site may be arbitrary: only this parameterized family
is netted, not the space of all Lipschitz functions.

\emph{Lower-degree induction and restoration.}
Apply induction to $H/H_r\to H_F/H_{F,r}$, with marked polynomial
$P_H\bmod H_{F,r}$ and the external test just constructed.  Its
fixed polynomial output lifts jointly with the exact map $P_H$ to
one weighted polynomial $h^U\in H$ by the rational filtered fiber
isomorphism.  Then $E^Uh^UR^U$ works for the averaged test and
projects to $P_F^*$.  Multiply it on the left by
$\rho(P_F(P_F^*)^{-1})$ to restore the formal marked projection
$P_F$; this correction is the identity on every retained slice.
Project to $G^0/L$ and lift polynomially to $G^0$ by a filtered
logarithmic section.  The original marked projection is unchanged
because $L\subset\ker\pi$.

Finally, if a path has averaged score at least $\delta$ and
$|a|\le B$, the set of $l\in L/(L\cap\Gamma^0)$ giving original
score at least $\delta/2$ has measure at least $\delta/(2B)$.
Otherwise the integral would be smaller than $\delta$, by bounding
the score by $B$ on that set and $\delta/2$ off it.  Fubini supplies
one fixed $l$ working on an early fraction of paths.  Its left
translation preserves the marked map and gives the desired output.

All dimensions count only polynomially many pivots.  Linear systems
at each grade are solved jointly.  Marked matching has at most $s$
grades and the degree induction at most $s$ steps.  Thus only a
bounded number of polynomial budget increases are used.  These
bounds are early and can all be fixed before the later scales.
\end{proof}

\begin{remark}[Uniformity in path-dependent tests]
\label{lifting:globalization-gap}
The proof uses the sampler's early-uniform testing seminorm for
simultaneous approximation of arbitrary masked inputs against all
path-dependent pivot tests.  Marked matching, exact reconstruction,
and the function-net bound are then deductions from this assertion
and Lemma~\ref{prelim:step-drop}.  These are the sampling and symbol
inputs to Proposition~\ref{lifting:candidate-globalization}.
\end{remark}

\begin{lemma}[A fixed patch evaluation on many paths]
\label{lifting:fixed-patch-function}
Suppose an early fraction of paths have positive scores of early
size against a real input of early bounded supremum, using patches of degree at most
$s$, rank at most $d_*$, and early kernel Lipschitz complexity.
The bound $d_*$ and all parameter dimensions are included in the
early structural budget.
After early losses in score and productive fraction, those patches
can be replaced by evaluations of one fixed Lipschitz function on
one fixed triangular patch group.  The local polynomial maps may
still vary with the path and have unrestricted coefficients.  The
fixed function has early complexity and at most $d_*$ patch slots.
\end{lemma}

\begin{proof}
First pigeonhole the rank and slot weights.  Let $D\le d_*$ be this
rank, $L\ge1$ a common upper bound on kernel Lipschitz constants, and
$\delta$ a common positive lower bound on score.
Normalize all residual polynomial coefficients modulo integers in
increasing slot order, using triangular integral changes of slot
variables.  This preserves the patch value, kernel, and weighted
degree.  Choose an accuracy
$\varepsilon\ll\delta/(1+\|a\|_\infty)$ and a grid mesh
$h\le\min(1/32,\varepsilon/L)$ depending only on early bounds.
Extend each kernel $\Phi$ by zero as required in the patch definition.
Tensor-product piecewise linear interpolation gives
\[
 \widetilde\Phi(x)=\sum_{v\in\mathcal V}\Phi(v)\,\Psi_h(x-v),
 \qquad 0\le\Phi(v)\le1,
 \qquad \|\widetilde\Phi-\Phi\|_\infty\le\varepsilon,
\]
where $\Psi_h$ is one nonnegative tensor-product tent of cap one,
supported in $[-h,h]^D$.  The vertices lie in a fixed slight
enlargement of the support and
$\log|\mathcal V|=O(D\log(2+L/\varepsilon))$.
The joint support lies in $[-1/4-h,1/4+h]^D$, of coordinate width
less than one.  Triangularity permits at most one integer slot tuple
there.  Substitution in the patch therefore changes its value by at
most $\varepsilon$, with no loss from counting tuples.  Its score
remains at least $\delta/2$.

All interpolation coefficients are nonnegative.  Some translated
tent patch consequently has score at least
$\delta/(2|\mathcal V|)$.  Absorb its center into the constant terms
of the residual polynomials; every path now uses the same kernel
$\Psi_h$.  These constant shifts leave the homogeneous
weight-preserving coefficients unchanged.

Write its triangular change as $S^0Y$, with $S^0$ the homogeneous
weight-preserving part and $Y$ in the fixed strictly lowering group.
All coefficients of $S^0$ are uniformly bounded and their number is
polynomial in $D$ for fixed $s$.  On a bounded fundamental region for
$Y$, triangular inversion gives a uniform early bound for all
integral tuples that can contribute for such an $S^0$ or a nearby
coefficient mesh point.  There are an early number of such tuples.
Polynomial substitution and the Lipschitz bound of $\Psi_h$ give an
early uniform continuity bound for
\[
          S^0\longmapsto
          \sum_{b\in\Z^D}\Psi_h(S^0Y(b))
\]
in supremum norm on the patch nilmanifold.  Mesh these finitely many
coefficients sufficiently finely, preserving triangular homogeneous
form.  The logarithmic mesh size is polynomial in the early budget.
A mesh replacement loses less than half the retained score, and
pigeonholing its index gives one fixed function on an early fraction
of paths.  No coefficients of the polynomial maps $Y$ were
pigeonholed.
\end{proof}

\begin{corollary}[Globalizing returned patches]
\label{lifting:globalize-patches}
Under the hypotheses of Proposition~\ref{lifting:candidate-globalization}
and sampling Property~\ref{sampling:property-transfer}, let
$0\le f\le1$ and $0\le\lambda\le1$.
A family as in Lemma~\ref{lifting:fixed-patch-function}, with real
input $a=f-\lambda$, yields an ambient patch of
rank at most $d_*+\sum_h\dim W_h$, with positive score after
the prescribed fixed multiplicative discount of the target density.
Its logarithmic inverse score, output slice costs, and complexity
may include the subspace heights and later parameters, as allowed
in sampling Property~\ref{sampling:property-transfer}.  They have polynomial bounds in that
augmented log budget, rather than necessarily early bounds.
\end{corollary}

\begin{proof}
Apply Lemma~\ref{lifting:fixed-patch-function} and then
Proposition~\ref{lifting:candidate-globalization} with trivial
initial marked factor.  Its fixed weighted polynomial remains in
the same triangular group, so its evaluation is a fixed patch
formula with $d_*$ slots and formal lift arguments.
Sampling Property~\ref{sampling:property-transfer} introduces at most $\sum_h\dim W_h$ slots to
determine those arguments, with the asserted score and discount.
\end{proof}

\subsection{Removing the lowest patch layer}

The removal of an old layer has an elementary part that can be proved
without any sampling theorem.

\begin{lemma}[Freezing a lowest layer]
\label{lifting:freeze-layer}
Consider a triangular patch with $D$ lowest-weight slots, all of
weight $j$.  Write their residual as
\[
                  r(t,b)=(I-M)(b-P(t)),
\]
where $M$ is strictly lower triangular.  Suppose $I(t)$ is an
integer-valued polynomial of ordinary degree at most $j$ on the
integer parameter box, and
\[
 \sup_t\left\|(I-M)\bigl(P(t)-I(t)-P(t_*)+I(t_*)\bigr)\right\|_\infty
       \le\varepsilon<\tfrac1{12}.
\]
Suppose the original patch has a contributing tuple at $t_*$, whose
first $D$ entries are $b_*$.  Then the integer polynomial
\[
                        b^*(t)=b_*+I(t)-I(t_*)
\]
is the only possible choice of its first $D$ slots anywhere that the
original patch contributes.  Substituting it into the higher slots
and freezing the first residual at $r(t_*,b_*)$ produces a patch with
these $D$ slots removed.  If the original kernel has Lipschitz
constant $L_\Phi$ for the maximum norm, the two patch values differ
pointwise by at most $L_\Phi\varepsilon$.
\end{lemma}

\begin{proof}
The displayed oscillation bound shows that the residual of $b^*(t)$
lies within $\varepsilon$ of $r(t_*,b_*)$, and hence lies in
$[-1/3,1/3]^D$.  If a different integer tuple $b$ first differs in
coordinate $i$, triangularity makes the difference between its
$i$th residual and that of $b^*(t)$ a nonzero integer.  The former
cannot lie in $[-1/4,1/4]$, because $1/3+1/4<1$.
Thus every contributing tuple has first block $b^*(t)$.

For this fixed block the higher residual formulas are identical in
the original and frozen tests.  The union of higher tuples contributing
to either kernel has size at most one, by triangularity.  For any such tuple the Lipschitz
bound gives an error at most $L_\Phi\varepsilon$.  The same argument
holds when one kernel value is zero, since the kernel is a globally
Lipschitz function supported in the prescribed box.  Summing over
the possible higher tuples therefore gives the pointwise bound.
Finally, substitution of a degree-at-most-$j$ polynomial for a
weight-$j$ variable preserves the weighted degree bounds of every
higher residual.
\end{proof}

\begin{lemma}[Rank-preparation bookkeeping]
\label{lifting:rank-potential}
Suppose preparation of $D$ coordinates of weight $j$ proceeds by
the following operation.  At some tag $h\ge1$, replace a space
$W_h$ by a codimension-one subspace and, if $h>1$, place one
one-dimensional contribution in a new tagged copy at $h-1$.
At $h=1$ absorb that contribution into degree zero.  Starting with
$W_j=\R^D$ and no other positive tags, at most $jD$ such operations
are possible.  At every stage
\[
                         \sum_h\dim W_h\le D.
\]
Only polynomially many ambient copy coordinates are needed when the
degree is fixed.
\end{lemma}

\begin{proof}
There are two dimension counts.  The unweighted sum
$\sum_h\dim W_h$ starts at $D$.  Each operation removes one
dimension and adds at most one, so this sum never increases.
The weighted sum $\sum_h h\dim W_h$ starts at $jD$ and decreases
by at least one at every operation, giving the bound on their number.

These sums count the dimensions of the spaces in which the polynomials
take values, not the numbers of their ambient coordinates.  A new copy
may have many ambient coordinates, but its incoming polynomial is a
scalar multiple of one removed vector and has a one-dimensional value
space.  There are at most $jD$ copies.  Each can be placed in a direct
sum of the original coordinates and copies from higher tags.  The
descent has depth at most $j\le s$, so the total ambient dimension
is polynomial in $D$ for fixed $s$.
\end{proof}

\begin{lemma}[Preparation of an old layer]
\label{lifting:prepare-layer}
Let $D\ge1$ and let $P:\R^n\to\R^D$ have degree at most $1\le j\le s$, with arbitrary
real coefficients, and prescribe a rank threshold $R\ge2$ and
oscillation tolerance $\delta>0$.  On sufficiently large parent
sides, there is a parameter slice $u=\iota(v)$ on which
\[
             P\circ\iota=c+I+E+\sum_{h\le j}S_h C_h.
\]
Here $c$ is constant, $I$ has integral nonconstant coefficients,
$E$ has oscillation at most $\delta$ in maximum norm,
$C_h$ has degree at most $h$, takes values in a rational space
$W_h\subseteq V_h$, and is high rank to threshold $R$ in the final
indexing.  The ambient space at tag $h$ is a direct sum
$V_h=\bigoplus_{\nu\in\mathcal I_h}\R^D$, with $\mathcal I_h$
indexing tagged copies of the original coordinate space, and
\[
 S_h:V_h\longrightarrow\R^D,\qquad
 S_h((z_\nu)_{\nu\in\mathcal I_h})
       =\sum_{\nu\in\mathcal I_h}z_\nu.
\]
Thus $\|S_h\|_{\infty\to\infty}\le|\mathcal I_h|$, independently
of the subspace heights.  Moreover
$\sum_h\dim W_h\le D$; all ambient copy dimensions are bounded
by a polynomial in $n+D$ for fixed $s$.  Rational log heights of the
spaces $W_h$ and auxiliary sections, and log slice costs, are polynomial in
$n+D+\log R+\log(2+\delta^{-1})$.
The slice may be selected to preserve the normalized average of any
specified real score function on the parent box.
\end{lemma}

\begin{proof}
Start with $C_j=P$ and $W_j=\R^D$.  If high rank fails at tag $h$,
there is an integer row $\lambda$ of norm at most $R$, nonzero on
$W_h$, for which the top of $\lambda C_h$ is rational plus slow.
Choose a rational $e\in W_h$ with $\lambda e=1$, and replace $C_h$
by $C_h-e\lambda C_h$, whose values lie in
$W_h\cap\ker\lambda$.  Write the removed term as its rational
top, its slow top, and a polynomial of degree at most $h-1$.

Restrict to residue classes clearing the denominators of the
rational top, including those of $e$.  Every positive-degree
coefficient then becomes integral; the constant is absorbed in $c$.
Subdivide the normalized box so finely that the slow top oscillates
by at most a prescribed share of $\delta$, and absorb its value at
one point into $c$.  Move the lower-degree term, which takes values
in the line spanned by $e$, into a new copy at tag $h-1$, or absorb
it as a constant if $h=1$.  Use the actual unmodified score to select
a cell of the residue and subbox partition whose normalized average
is at least the preceding average.  This exists by averaging.
Equal or almost equal subdivisions keep every side at a controlled
relative size once the parent cutoffs are imposed.

After reindexing, repeat the rank tests at the same $R$, including
tags tested before a subsequent restriction.  Lemma~\ref{lifting:rank-potential}
shows that at most $jD$ failures can be removed.  At termination every
remaining tag is high rank in the final indexing.  Previously
integral nonconstant coefficients remain integral under integer
affine reindexing.  Previously allocated oscillation errors can only
decrease on restriction.  Thus allocating $\delta/(jD)$ per step,
with the polynomial norm bound of the copy-summing maps included,
gives the asserted total error.

We verify that the rational estimates do not undergo one polynomial
exponentiation per cut.  Fix a polynomial bound $m$ for all copy
coordinates.  At tag $h$ let $B_h$ be the block-diagonal basis matrix
of its original and incoming spaces, before cuts at that tag, and
stack all imposed rows into $A_h$.  The current space is
$B_h\ker(A_hB_h)$.  Extend earlier rows by zeros whenever a copy
arrives.  To choose the next section, solve jointly
\[
 A_hB_hz=0,\qquad \lambda B_hz=1,
 \qquad e=B_hz.
\]
The system is consistent because $\lambda$ is nonzero on the
current space.  Clearing the polynomial number of denominators and
using one nonzero minor gives
\[
 \log\operatorname{ht}(e),\ \log\operatorname{ht}(W_h)
 \le m^{C_s}\bigl(1+\log R+\log\operatorname{ht}(B_h)\bigr).
\]
Crucially, $B_h$ contains sections coming from higher tags, not
sections previously chosen at this same tag.  Thus the displayed
polynomial recurrence is composed at most $s$ times.  It bounds both
$\|e\|_\infty$ and a common denominator $d_e$ of its coordinates.

For completeness, if $q\le R$ clears the rational scalar top, the
modulus $a=qd_e$ clears the nonconstant coefficients of its product
with $e$ on every residue class $u=r+av$: every positive-degree
term in the expansion contains a factor $a$.  Its arbitrary
numerators do not affect this assertion.  In normalized variables
the slow top has at most $\binom{n+s}{s}$ monomials with coefficients
at most $R$.  If $K$ bounds the copy-summing maps, subdivision into
$B_0$ pieces per coordinate costs an oscillation at most
\[
 \frac{Ks\binom{n+s}{s}R\|e\|_\infty}{B_0}.
\]
Choose $B_0$ large enough to make this at most $\delta/(jD)$.
Then $\log(2aB_0)$ is polynomial in the asserted budget.  On
sufficiently large current sides, nearly equal subdivisions in
each residue class give every cell side at least $H_i/(2aB_0)$.
Consequently each cut has that polynomial log cost, and summing
over at most $jD$ cuts gives both the claimed cost and a uniform
sufficient parent-side cutoff.  Copy dimensions are bounded by a
fixed polynomial majorant of $D(1+jD)^j$, by the bounded descent
depth, independently of orientations and heights.
\end{proof}

After successful preparation and sampling, the integer polynomial
$I(t)$ in Lemma~\ref{lifting:freeze-layer} is the prepared
integral-coefficient contribution plus
$\sum_h S_h(\beta_h(\psi(t)))$.  The small residual oscillation is the sum of the
preparation error and the chart widths.  Choosing that error small
relative to the input absolute score permits the recursive call with
the original density threshold, losing only absolute surplus.
This choice must be made before the rank threshold and subspace
heights become available.

\begin{proposition}[Linear rank accounting]
\label{lifting:linear-rank}
Assume the following three assertions, with their permitted
early or augmented logarithmic cost bounds:
\begin{enumerate}
\item rank preparation and conditioned sampling remove the lowest
      $D$ slots of weight $j$, retaining positive absolute score and
      producing spaces with $\sum_h\dim W_h\le D$;
\item a family of recursively returned patches of rank at most $d_*$
      can be globalized in the same patch group as in
      Proposition~\ref{lifting:candidate-globalization};
\item transferring this fixed global patch to the parent costs at
      most $\sum_h\dim W_h$ new slots, and uses an arbitrarily
      prescribed fixed multiplicative discount of the target density.
\end{enumerate}
Assume also the unconditioned base case returns rank at most $d_0$
and permits an arbitrarily prescribed fixed density discount.
Then processing an old patch with $d$ slots returns rank at most
$d_0+d$.  For fixed $s$, its total multiplicative density discount
can be made at most any prescribed $\epsilon_0>0$.
\end{proposition}

\begin{proof}
Let $D_j$ be the number of original slots of weight $j$.
After the lowest remaining layer $j$ is removed, the recursive call
acts on $d-D_j$ old slots.  On returning, the old test is discarded:
only the returned patch is globalized and charged the new
$\sum_h\dim W_h$ slots.  Consequently the rank satisfies
\[
 R(D_1,\ldots,D_s)
       \le R(0,\ldots,0,D_{j+1},\ldots,D_s)+D_j,
 \qquad R(0,\ldots,0)\le d_0.
\]
Iterating gives
$R\le d_0+\sum_jD_j=d_0+d$.
There are at most $s$ upward passages and one base passage.  Taking
each fixed discount $\delta\le\epsilon_0/(s+1)$ makes their product
at least $1-(s+1)\delta\ge1-\epsilon_0$.
\end{proof}

\begin{remark}[Order of parameters in the relative use]
\label{lifting:relative-gap}
The recursive input budget is fixed from the old
patch, the dimensions of all possible copies, the desired chart
widths, and the retained absolute score.  It must not include
$\log R$, $\log L$, or the eventual heights of the prepared spaces.
The sampler dimension must be fixed from structural dimensions
before these recursive accuracy requests; a mere polynomial bound
on that dimension allowing dependence on later requested precisions
would create a circular requirement.  This structural independence
is explicitly part of Proposition~\ref{sampling:constrained}.

The removed directions $e$ are stored in the unrestricted
coefficients of the copied polynomials.  The maps $S_h$ only add
coordinate copies, and their norms depend on copy counts, not on
the rational heights.  Therefore the widths needed in
Lemma~\ref{lifting:freeze-layer}, and the complexity of the remaining
kernel after freezing, can be prescribed from the structural budget.
The coefficients of the substituted higher residuals remain free.
Compute the recursive output, accuracy, and globalization budgets
with these data first.  Only then choose perturbations, the long-scale
range and $R$, perform Lemma~\ref{lifting:prepare-layer}, choose the
allowed long scale for its final spaces, and impose parent side
cutoffs including its late height and slice costs.  At most $s$ old
layers are processed, so this forecasting uses bounded depth.
This is the order used in the relative induction below, with sampling
supplied by Proposition~\ref{sampling:constrained} and step-drop by
Lemma~\ref{prelim:step-drop}.
\end{remark}

\begin{theorem}[Relative lifting from the sampling assertions]
\label{lifting:relative-theorem}
Assume Lemmas~\ref{prelim:step-drop} and~\ref{prelim:separation},
Proposition~\ref{sampling:constrained}, and the scalar transfer of
Proposition~\ref{sampling:scalar-transfer}, including their stated
uniform polynomial logarithmic bounds.  Then
Proposition~\ref{counting:relative-target} holds, with rank at most
$d_0+d$ in its increment alternative.  Its exponent $E$ depends
only on $s$, the fixed base dimension $n_0$, its fixed side-ratio
bound, and the prescribed discount $\varepsilon_0$.
\end{theorem}

\begin{proof}
Induct on the number of distinct positive weights in the old patch.
Fix a discount $\tau\le\varepsilon_0/(s+1)$ for each upward passage
and the base passage, with $0<\tau\le1/2$.  For early bookkeeping
replace the absolute rank bound by $\min(d_0,p)$: the rule's
complexity bound already bounds every actual returned rank by $p$.
The eventual rank estimate with the original $d_0$ only becomes
weaker.  We use the absolute rule solely through its
given dimension, cutoff, score, slice, rank, and complexity bounds;
no constant below depends on the particular rule.

\emph{Base case.}
If the old patch has no slots, it is a constant in $[0,1]$.
The positive old score gives $\E f\ge a+e^{-p}$: if the constant
old patch is $c$, then $c>0$ and $\E f-a\ge e^{-p}/c\ge e^{-p}$.
Use the unconditioned sampler in dimension exactly $n_0$, with
distinct comparable prime sides beyond the absolute cutoff.  Its
spatial law is dominated by the product-box law by an early factor,
as specified in the unconditioned variant of
Proposition~\ref{sampling:constrained}; hence scalar transfer applies
to this same ensemble.
Put $\sigma=e^{-p}$ and prescribe both the one-site error and the
noninjective probability to be at most $\sigma/16$.  If $Z$ is the
sample mean minus $a$, then $|Z|\le1$ and
$\E[Z1_{\mathrm{injective}}]\ge7\sigma/8$.  Since $Z\le1$,
the fraction of injective samples with $Z\ge\sigma/2$ is at least
$\sigma/4$.  These samples therefore have mean at least
$a+e^{-p}/2$.  The structural side-ratio constant in the absolute
rule is fixed with $R\ge2^{n_0+1}$.  Choose $L$ divisible by
$2^{n_0+1}$ and apply Bertrand's postulate in $n_0$ successive
disjoint dyadic intervals inside $[L/R,L]$.  This gives distinct
prime sides comparable to $L$.  Taking $L/R>e^p$ puts every side
beyond the rule's cutoff; these are all early requirements on $L$.

The pullback of $f$ on an injective affine sample is progression-free:
a nonzero vector progression in its box has two distinct points,
whose images remain distinct, and its image is an arithmetic
progression in the parent.  If the absolute rule is an impossibility
rule, these samples are already a contradiction.  Otherwise that
rule supplies a patch of rank at most $d_0$ and early positive score
against the pullback minus $\Lambda$ on a permitted sample slice.
Lemma~\ref{lifting:fixed-patch-function} and the unconditioned case
of Proposition~\ref{lifting:candidate-globalization} replace these
by one fixed ambient patch evaluation, still with at most $d_0$
slots, on an early fraction of samples.

Apply scalar transfer to this fixed ordinary niltest $g$, with
$\lambda'=(1-\tau)\Lambda$.  It has the required positive lower
bound, after an early budget increase.  Choose its fixed slack
$\epsilon=\tau/(2(1-\tau))$, so that
$(1+\epsilon)\lambda'=(1-\tau/2)\Lambda\le\Lambda$.
If every permitted ambient slice had score at most the prescribed
small error against $f-\lambda'$, scalar transfer would bound the
probability of the observed local positive scores by a smaller
early fraction.  Here is the finite union bound for the path-dependent
slices.  Let $b$ bound their total early costs, let $\delta$ be their
retained score and let $\rho$ be the productive fraction.  Each
original side is comparable to $L$, and each restricted side has
at least $Le^{-O(b)}$ points.  Once this is at least two, its common
modulus can be represented by an integer at most $e^{O(b)}$.
Enumerate these moduli and residues, and mesh the interval endpoints
in units of $L$ at resolution $\delta e^{-O(b)}/C_{n_0}$.  The change
in the normalized slice measure is at most $\delta/4$ in $\ell^1$
distance, after $L$ absorbs the integer-rounding error.  The score
integrand is bounded by one, so this is also a score bound.  There
are $N\le\exp(\mathcal P(b+\log\delta^{-1}))$ choices, independently
of $L$.  Choose the scalar log budget $p_{\rm sc}$ in advance so
that its complexity, modulus, side-length and domination hypotheses
hold and
\[
 e^{-p_{\rm sc}}<\min\{\delta/4,\rho/(2N)\}.
\]
Apply Proposition~\ref{sampling:scalar-transfer} to each fixed mesh
choice with its own uniform $t$ law.  Only the outer law of $(x,V)$
is changed to the actual dominated unconditioned law; no
path-dependent reweighting of $t$ is used.  If every permitted
ambient slice had score at most the resulting $e^{-P}$, the union
of the local positive-score events would have probability at most
$Ne^{-p_{\rm sc}}<\rho/2$, a contradiction.  All endpoint,
precision and domination budgets were fixed before choosing $L$.
The contradiction yields an ambient slice with early positive score
against $f-(1-\tau)\Lambda$, with the same $d_0$ slots.  This proves
the base increment case.

\emph{Removing the lowest weight.}
Suppose the lowest weight is $j$ and it has $D$ slots.  Normalize the
patch coefficients by the triangular integral substitutions in
Lemma~\ref{lifting:fixed-patch-function}.  The kernel is unchanged.
At this lowest weight, previous-slot terms are linear with constant
coefficients, now of absolute value at most one.  Thus its residual
block is $(I-M)(b-P(u))$, with
$\|I-M\|_{\infty\to\infty}\le1+D$ and $\deg P\le j$.
The coefficients of $P$ need not be bounded.
Forecast the polynomial number of all possible copy coordinates in
Lemma~\ref{lifting:prepare-layer}, and let $K$ bound
$\sum_h\|S_h\|_{\infty\to\infty}$ for every such configuration.
Put $\sigma=e^{-p}$ and prescribe the preparation error $\delta$
and common chart radius $w$ so that
\[
 (1+D)(\delta+2Kw)
 \le \min\{1/24,\sigma/(16\max(1,L_\Phi))\}.
\]
Only the early kernel, score and structural data enter these
choices, not the heights of the eventually prepared spaces.

The constrained sampler dimension is fixed structurally.  Consequently
the inductive input dimensions, kernel complexity, desired inverse
score, and the original absolute rule's bounds can all be put in an
early polynomial budget before choosing its testing accuracies or
long scales.  Compute the induction output and globalization budgets
first.  Then choose the later parameters in the sampler's required
order, choose $R$ for the entire allowed long-scale range, and
perform Lemma~\ref{lifting:prepare-layer}.  The initial side cutoff
includes its polynomial late height and slice costs.  The final
prepared spaces admit one of the allowed long scales, and their
high rank holds at the threshold already chosen for the whole range.

Preparation preserves the actual old score, which is at least
$\sigma$.  Prescribe the mixed-center one-site error and the bad-path
probability to be at most $\sigma/16$ each.  For the sampled old
score $Z$, $|Z|\le1$ and
$\E[Z1_{\rm good}]\ge7\sigma/8$.  Therefore at least a
$\sigma/4$ fraction of paths are good and have $Z\ge\sigma/2$.
Averaging in the centers gives one fixed center tuple with at least
that productive fraction.  On each such path
choose a point where the old patch contributes.  The prepared
integral part plus $\sum_h S_h(\beta_h(\psi(t)))$ is the integer polynomial
$I(t)$ of Lemma~\ref{lifting:freeze-layer}: the prepared integral
part stays integral under the affine integer map, and each
$\beta_h\circ\psi$ has integer coefficients and degree at most
$h\le j$.  Constant centers cancel in differences.  The residual
oscillation is at most $(1+D)(\delta+2Kw)$, so freezing loses at
most $\sigma/16$ in score.  The remaining patch has $d-D$ slots
and score at least $7\sigma/16$ against $f-a$.  Its kernel is
$\Phi(r_*,\cdot)$, with the same support and no larger Lipschitz
constant.  Its residual coefficients may be large, as permitted
in the patch definition.

Let $p_{\rm child}$ include the structural dimension, frozen kernel
complexity, $p$, and $\log(16/(7\sigma))$, and let $H_{\rm child}$
be the already computed inductive side cutoff for that budget.
Choose $L>H_{\rm child}$.  Freeze axes with $T_i<H_{\rm child}$
on a favorable fiber.  This preserves normalized score, and each
fixed axis costs at most $\log H_{\rm child}$; their total cost is
early.  Omit these singleton coordinates for induction, then
reinsert their fixed values when recording its output as a slice
of $\mathcal T$.  At least one free axis remains because the
sampler has a side of length $L$, and every free side is at least
$H_{\rm child}$.  The same absolute rule is valid with budget
$p_{\rm child}\ge p$: its cutoff, slice, complexity and inverse-score
bounds only become weaker, and $a\ge e^{-p_{\rm child}}$.
Apply induction to this positive-dimensional box and patch.  In the impossibility
case induction contradicts the retained score, so that alternative
is proved at every depth.  In the increment case it returns patches
and slices of early complexity and positive score at the inductively
discounted target, with rank at most $d_0+d-D$.

Discard the old test and globalize precisely these returned patches
by Corollary~\ref{lifting:globalize-patches}.  This costs at most
$\sum_h\dim W_h\le D$ slots and one additional fixed discount
$\tau$.  Therefore the rank is at most
\[
                  d_0+(d-D)+D=d_0+d.
\]
There are at most $s$ old-weight passages plus the base passage, so
the product of the target factors is at least
$(1-\tau)^{s+1}\ge1-\varepsilon_0$.  Lowering the target to the
one requested only increases its positive score.

Here is an explicit form of the final budget recurrence.  Let
$F_q(p)$ majorize all required logarithmic bounds when at most $q$
distinct old weights remain, including the sufficient parent log
side.  The base construction gives a polynomial $F_0$.  Structural
dimensions, the frozen kernel and its retained score admit
$p_{\rm child}\le(2+p)^{c_s}$, independently of $R,L$ and heights.
After computing the child output, all globalization and testing
requirements fit in
\[
 b_*\le\bigl(2+p+F_{q-1}((2+p)^{c_s})\bigr)^{c_s}.
\]
The quantitative sampler hypotheses choose perturbation, $A$ and
$R$ with logarithms bounded by a fixed polynomial in $b_*$, valid
throughout $A\le L\le A^{C_s(1+\sum_hd_h)}$.  Preparation's joint
height estimate then bounds its heights and slice costs by a
polynomial in $p+\log R$.  The allowed augmented output costs,
and the parent cutoff absorbing every restriction, are consequently
bounded by one more fixed polynomial in $b_*$.  Increasing a fixed
exponent if necessary gives
\[
 F_q(p)\le
 \bigl(2+p+F_{q-1}((2+p)^{c_s})\bigr)^{C_s}.
\]
The polynomial number of preparation cuts adds costs; it does not
add compositions to this recurrence.  Since $q\le s$, there is
one exponent $E$, depending only on the fixed parameters, bounding
all sufficient initial side cutoffs in logarithm, final complexities,
log inverse scores and slice costs by $(2+p)^E$.  This is the full uniform assertion of
Proposition~\ref{counting:relative-target}.
\end{proof}
